\section{What is at stake? Three policy levers}\label{sec:stake}

The introduction set out two families of policy: if the invention is the problem, invest in research capability; if
the inventor is, relieve owners' constraints through credit, fee relief or markets for technology. The renewal model
can price the levers that the evidence points to on a common scale. I simulate each patent on its own fee path to
expiry with the deterministic model of Section~\ref{sec:values}, change one thing at a time, and report the change
in the returns owners earn from their patents per filing-year cohort, in 2024 PPP dollars and as a share of national
R\&D spending, for China, Korea and Brazil (Table~\ref{tab:levers}).\footnote{\label{fn:levers}China, Korea and Brazil are the three emerging-market
offices with fee discounts for small owners, which the first lever needs. Korea records no sales of granted patents
and Brazil's records, which name the new owner, show too few sales of granted patents by domestic owners to unrelated buyers (8 at the
ages studied), so markets for technology are priced for China only. The Czech
Republic records no ownership transfers and has no fee discount for small owners, and the Russian values are not
identified (Section~\ref{sec:values}). Patents a year: in China 4{,}100 of
small firms, 1{,}800 of mid-sized firms, 5{,}200 of large firms and 6{,}600 of universities; in Korea 4{,}700, 3{,}800,
26{,}900 and 3{,}300; in Brazil 45, 20, 18 and 43. Firm size: inventions filed worldwide in 2000--2005 by the owner. National R\&D: gross domestic expenditure on
R\&D, 2000--2005 average \citep{WDI}. Decay factor $d = 0.85$; fees after 2024 at 2024 prices. Fee relief: every eligible
owner (firms with fewer than 100 inventions, universities) pays the discounted fees; in Brazil 40\% of the
standard annuity (50\% from August 2025), simulated for patents granted from June 2009, whose fees are known. Markets for technology:
returns multiplied, from one year after grant, by a factor calibrated so that the model's survival five years after a
sale, 1--5 years after grant, rises by the extra post-sale survival of sales to unrelated buyers in China relative to
Germany and France, counted from grant: 3.8 points for small firms and 4.8 for mid-sized firms,
giving factors of 1.22 and 1.40; none for large firms. For universities the factor, 3.5, matches Chinese
universities' own premium of 46 points less the 7 points in Germany and France. University patenting: mean log returns of Chinese
universities raised to those of foreign universities in China (upper bound) or to the foreign level plus the average
domestic--foreign university gap in the United States, Germany, the United Kingdom and France (lower bound).} Fees saved are transfers from the state to owners,
so they are not counted in the table.

\paragraph{Fee relief.} Suppose every eligible owner (firms with fewer than 100 inventions and universities) took up
the Chinese, Korean and Brazilian fee discounts. Owners of a Chinese cohort would pay \$8.0 million less in fees,
owners of a Korean one \$1.8 million less and owners of a Brazilian one \$0.2 million less, but they would keep their
patents only marginally longer, and the returns they earn would rise by about half a million dollars in China, 1.3
million in Korea and 0.1 million in Brazil. Fees are small relative to what patents are worth, so cutting them
changes few renewal decisions. Brazil is the partial exception: its patents are worth little relative to its fees,
and the discount would raise the value of eligible owners' patents by 37--50\%, but they hold only about 110 patents
a cohort.

\paragraph{Markets for technology.} Suppose that, from one year after grant, domestic owners' patents earned what they
would earn in the hands of a buyer. For each group I calibrate the buyer's advantage to the extra survival of patents sold to
unrelated buyers in Section~\ref{sec:constraints}, relative to sellers of the same size in Germany and France,
counted from grant to match the model. In China this requires returns about 20\% higher for small firms and 40\% for
mid-sized firms, and the returns from a cohort of small firms' patents would rise by about \$130 million, or \$280
million with mid-sized firms. Both targets are imprecise, so these are orders of magnitude; large firms show no
premium. Korea records no sales of granted patents, so the lever is not priced there. For Chinese universities, whose
sold patents survive 45 points longer, the calibrated advantage is a 3.5-fold increase in returns and the gain
\$460 million. Each figure is an upper bound, since it reallocates every patent of the group.

\paragraph{University patenting.} Suppose instead that Chinese universities' patents earned what foreign
universities' patents earn in China at equal citations. Each cohort would be worth about \$320 million more; with
the typical university home premium of the advanced-economy offices as the benchmark, about \$180 million. This and
the university row under markets for technology are two routes to the same lost value, improving what universities
patent or moving their patents to owners who use them, and should not be added.

\begin{table}[t]
\centering
\begin{threeparttable}
\caption{Returns at stake per filing-year cohort, by lever and owner group (millions of 2024 PPP dollars)}\label{tab:levers}
\small
% Generated by em-patent-value code/38_paper_exhibits.py (policy_levers.csv, step 26); do not edit by hand
\begin{tabular}{lccccc}
\toprule
 & \multicolumn{3}{c}{Fee relief} & Markets for technology & University patenting \\
\cmidrule(lr){2-4}\cmidrule(lr){5-5}\cmidrule(lr){6-6}
 & China & Korea & Brazil & China & China \\
\midrule
Small firms (1--9)          & 0.07 & 0.49 & 0.05 & 133 &  \\
Mid-sized firms (10--99)    & 0.02 & 0.42 & 0.03 & 146 &  \\
Large firms (100+)          & n.e. & n.e. & n.e. & 0.00 &  \\
Universities                & 0.43 & 0.34 & 0.04 & 462 & 178--319 \\
\addlinespace
All owners                  & 0.52 & 1.3 & 0.11 & 742 & 178--319 \\
\quad \% of national R\&D & 0.00 & 0.00 & 0.00 & 0.79 & 0.19--0.34 \\
\bottomrule
\end{tabular}
\begin{tablenotes}\footnotesize
\item \emph{Notes:} Change in the returns domestic owners earn from their patents, excluding fees, for one filing-year cohort; deterministic renewal model, each patent simulated on its own fee path. Markets for technology: upper bounds, calibrated to the extra post-sale survival of arm's-length sales; the firm targets are imprecise (\SuppTab{S5}{tab:sales}). The two university entries are alternative routes to the same value and are not additive. n.e.: not eligible. Details in footnote~\ref{fn:levers}.
\end{tablenotes}
\end{threeparttable}
\end{table}


\paragraph{Reading the numbers.} Two cautions apply. The figures are private returns to owners, not welfare: they
leave out knowledge spillovers and the gain to users when a patent enters the public domain. The levels are rough: the
decay rate is fixed rather than estimated in China, the sale targets are imprecise, and in the learning model, where
owners keep low-value patents for the option of later renewals, the same survival gains require far larger changes in
returns and give levels several times higher that are not identified (Supplemental Appendix Table~S9). The ordering survives both: relieving owners of fees is worth
about a million dollars a cohort or less, while helping patents reach owners who can use them and improving what Chinese
universities patent are worth hundreds of millions, two to three orders of magnitude more.
